% PDF page 1
\source{gilbarg-trudinger:page-0001}

% Cover page transcription.

CLASSICS IN MATHEMATICS

David Gilbarg \(\cdot\) Neil S. Trudinger

Elliptic
Partial Differential
Equations
of Second Order

Springer

% PDF page 2
\source{gilbarg-trudinger:page-0002}
Classics in Mathematics

David Gilbarg $\cdot$ Neil S. Trudinger \hfill Elliptic Partial Differential Equations
of Second Order

% PDF page 3
\source{gilbarg-trudinger:page-0003}

Springer

Berlin
Heidelberg
New York
Barcelona
Hong Kong
London
Milan
Paris
Singapore
Tokyo

% PDF page 4
\source{gilbarg-trudinger:page-0004}

David Gilbarg was born in New York in 1918, and was educated there through undergraduate school. He received his Ph.D. degree at Indiana University in 1941. His work in fluid dynamics during the war years motivated much of his later research on flows with free boundaries.

He was on the Mathematics Faculty at Indiana University from 1946 to 1957 and at Stanford University from 1957 on. His principal interests and contributions have been in mathematical fluid dynamics and the theory of elliptic partial differential equations.

Neil S. Trudinger was born in Ballarat, Australia in 1942. After schooling and undergraduate education in Australia, he completed his Ph.D. at Stanford University, USA in 1966. He has been a Professor of Mathematics at the Australian National University, Canberra since 1973. His research contributions, while largely focused on non-linear elliptic partial differential equations, have also spread into geometry, functional analysis and computational mathematics. Among honours received are Fellowships of the Australian Academy of Science and of the Royal Society of London.

% PDF page 5
\source{gilbarg-trudinger:page-0005}
\begin{center}
David Gilbarg \(\bullet\) Neil S. Trudinger

\vspace{2.5em}

\Large\textbf{Elliptic Partial}

\Large\textbf{Differential Equations}

\Large\textbf{of Second Order}

\vspace{2.5em}

Reprint of the 1998 Edition

\vfill

Springer
\end{center}

% PDF page 6
\source{gilbarg-trudinger:page-0006}
\begin{center}
\begin{tabular}{@{}p{0.42\textwidth}p{0.42\textwidth}@{}}
\textbf{David Gilbarg} &
\textbf{Neil S. Trudinger} \\
Stanford University &
The Australian National University \\
Department of Mathematics &
School of Mathematical Sciences \\
Stanford, CA 94305-2125 &
Canberra ACT 0200 \\
USA &
Australia \\
\textit{e-mail:} gilbarg@math.stanford.edu &
\textit{e-mail:} neil.trudinger@anu.edu.au
\end{tabular}
\end{center}

\vspace{2.5em}

\noindent\rule{\textwidth}{0.4pt}

\noindent\textbf{Originally published as Vol. 224 of the}\\
\textit{Grundlehren der mathematischen Wissenschaften}

\noindent\rule{\textwidth}{0.4pt}

\vspace{2.0em}

\noindent\textbf{Cataloging-in-Publication Data applied for}

\vspace{1.5em}

\noindent\textbf{Die Deutsche Bibliothek - CIP-Einheitsaufnahme}

\vspace{1.5em}

\noindent\textbf{Gilbarg, David:}\\
Elliptic partial differential equations of second order / David Gilbarg; Neil S. Trudinger.- Reprint of the 1998\\
ed.- Berlin; Heidelberg; New York; Barcelona; Hong Kong; London; Milan; Paris; Singapore; Tokyo:\\
Springer, 2001\\
(Classics in mathematics)

\vspace{3.0em}

\noindent\textbf{Mathematics Subject Classification (2000): 35Jxx}

\vspace{3.0em}

\noindent\textbf{ISSN 1431-0821}\\
\textbf{ISBN-13: 978-3-540-41160-4}\hfill \textbf{e-ISBN-13: 978-3-642-61798-0}\\
\textbf{DOI: 10.1007/978-3-642-61798-0}

\vspace{2.0em}

\noindent This work is subject to copyright. All rights are reserved, whether the whole or part of the material is
concerned, specifically the rights of translation, reprinting, reuse of illustrations, recitation, broadcasting,
reproduction on microfilm or in any other way, and storage in data banks. Duplication of this publication or
parts thereof is permitted only under the provisions of the German Copyright Law of September 9, 1965, in its
current version, and permission for use must always be obtained from Springer-Verlag. Violations are liable
for prosecution under the German Copyright Law.

\vspace{2.0em}

\noindent Springer-Verlag Berlin Heidelberg New York\\
a member of Springer Science+Business Media

\vspace{1.0em}

\noindent \textcopyright\ Springer-Verlag Berlin Heidelberg 2001

\vspace{2.0em}

\noindent The use of general descriptive names, registered names, trademarks etc. in this publication does not imply,
even in the absence of a specific statement, that such names are exempt from the relevant protective laws and
regulations and therefore free for general use.

\vspace{2.0em}

\noindent Printed on acid-free paper\hfill SPIN 10982136\hfill 41/3111ck-5 4 3 2 1

% PDF page 7
\source{gilbarg-trudinger:page-0007}

% Preface to the Revised Third Printing.

This revision of the 1983 second edition of ``Elliptic Partial Differential Equations
of Second Order'' corresponds to the Russian edition, published in 1989, in which
we essentially updated the previous version to 1984. The additional text relates to
the boundary H\"older derivative estimates of Nikolai Krylov, which provided a
fundamental component of the further development of the classical theory of
elliptic (and parabolic), fully nonlinear equations in higher dimensions. In
our presentation we adapted a simplification of Krylov's approach due to Luis
Caffarelli.

The theory of nonlinear second order elliptic equations has continued to
flourish during the last fifteen years and, in a brief epilogue to this volume, we
signal some of the major advances. Although a proper treatment would necessi-
tate at least another monograph, it is our hope that this book, most of whose text
is now more than twenty years old, can continue to serve as background for these
and future developments.

Since our first edition we have become indebted to numerous colleagues, all
over the globe. It was particularly pleasant in recent years to make and renew
friendships with our Russian colleagues, Olga Ladyzhenskaya, Nina Ural'tseva,
Nina Ivochkina, Nikolai Krylov and Mikhail Safonov, who have contributed so
much to this area. Sadly, we mourn the passing away in 1996 of Ennio De Giorgi,
whose brilliant discovery forty years ago opened the door to the higher-dimen-
sional nonlinear theory.

October 1997

David Gilbarg $\cdot$ Neil S. Trudinger

% PDF page 8
\source{gilbarg-trudinger:page-0008}

\noindent\textbf{Preface to the First Edition}

\medskip

\noindent This volume is intended as an essentially self-contained exposition of portions of the theory of second order quasilinear elliptic partial differential equations, with emphasis on the Dirichlet problem in bounded domains. It grew out of lecture notes for graduate courses by the authors at Stanford University, the final material extending well beyond the scope of these courses. By including preparatory chapters on topics such as potential theory and functional analysis, we have attempted to make the work accessible to a broad spectrum of readers. Above all, we hope the readers of this book will gain an appreciation of the multitude of ingenious barehanded techniques that have been developed in the study of elliptic equations and have become part of the repertoire of analysis.

\noindent Many individuals have assisted us during the evolution of this work over the past several years. In particular, we are grateful for the valuable discussions with L. M. Simon and his contributions in Sections 15.4 to 15.8; for the helpful comments and corrections of J. M. Cross, A. S. Geue, J. Nash, P. Trudinger and B. Turkington; for the contributions of G. Williams in Section 10.5 and of A. S. Geue in Section 10.6; and for the impeccably typed manuscript which resulted from the dedicated efforts of Isolde Field at Stanford and Anna Zalucki at Canberra. The research of the authors connected with this volume was supported in part by the National Science Foundation.

\medskip

\noindent August 1977 \hfill David Gilbarg \hfill Neil S. Trudinger

\noindent Stanford \hfill Canberra

\bigskip

\noindent\hrulefill

\noindent \textit{Note:} The Second Edition includes a new, additional Chapter 9. Consequently Chapters 10 and 15 referred to above have become Chapters 11 and 16.

% PDF page 9
\source{gilbarg-trudinger:page-0009}

Table of Contents

Chapter 1. Introduction . . . . . . . . . . . . . . . . . . . . . . 1

Part I. Linear Equations

Chapter 2. Laplace’s Equation . . . . . . . . . . . . . . . . . . 13
  2.1.  The Mean Value Inequalities . . . . . . . . . . . . . . . 13
  2.2.  Maximum and Minimum Principle . . . . . . . . . . . . . . 15
  2.3.  The Harnack Inequality . . . . . . . . . . . . . . . . . . 16
  2.4.  Green’s Representation . . . . . . . . . . . . . . . . . . 17
  2.5.  The Poisson Integral . . . . . . . . . . . . . . . . . . . 19
  2.6.  Convergence Theorems . . . . . . . . . . . . . . . . . . . 21
  2.7.  Interior Estimates of Derivatives . . . . . . . . . . . . . 22
  2.8.  The Dirichlet Problem; the Method of Subharmonic Functions . 23
  2.9.  Capacity . . . . . . . . . . . . . . . . . . . . . . . . . 27
Problems . . . . . . . . . . . . . . . . . . . . . . . . . . . . . 28

Chapter 3.  The Classical Maximum Principle . . . . . . . . . . . 31
  3.1.  The Weak Maximum Principle . . . . . . . . . . . . . . . . 32
  3.2.  The Strong Maximum Principle . . . . . . . . . . . . . . . 33
  3.3.  Apriori Bounds . . . . . . . . . . . . . . . . . . . . . . 36
  3.4.  Gradient Estimates for Poisson’s Equation . . . . . . . . 37
  3.5.  A Harnack Inequality . . . . . . . . . . . . . . . . . . . 41
  3.6.  Operators in Divergence Form . . . . . . . . . . . . . . . 45
Notes . . . . . . . . . . . . . . . . . . . . . . . . . . . . . . 46
Problems . . . . . . . . . . . . . . . . . . . . . . . . . . . . . 47

Chapter 4.  Poisson’s Equation and the Newtonian Potential . . . . 51
  4.1.  Hölder Continuity . . . . . . . . . . . . . . . . . . . . . 51
  4.2.  The Dirichlet Problem for Poisson’s Equation . . . . . . . 54
  4.3.  Hölder Estimates for the Second Derivatives . . . . . . . . 56
  4.4.  Estimates at the Boundary . . . . . . . . . . . . . . . . . 64

% PDF page 10
\source{gilbarg-trudinger:page-0010}
\begin{quote}
4.5. \ Hölder Estimates for the First Derivatives \dotfill 67

Notes \dotfill 70

Problems \dotfill 70

Chapter 5. \ Banach and Hilbert Spaces \dotfill 73

5.1. \ The Contraction Mapping Principle \dotfill 74
5.2. \ The Method of Continuity \dotfill 74
5.3. \ The Fredholm Alternative \dotfill 75
5.4. \ Dual Spaces and Adjoints \dotfill 79
5.5. \ Hilbert Spaces \dotfill 80
5.6. \ The Projection Theorem \dotfill 81
5.7. \ The Riesz Representation Theorem \dotfill 82
5.8. \ The Lax-Milgram Theorem \dotfill 83
5.9. \ The Fredholm Alternative in Hilbert Spaces \dotfill 83
5.10. \ Weak Compactness \dotfill 85
Notes \dotfill 85
Problems \dotfill 86

Chapter 6. \ Classical Solutions; the Schauder Approach \dotfill 87

6.1. \ The Schauder Interior Estimates \dotfill 89
6.2. \ Boundary and Global Estimates \dotfill 94
6.3. \ The Dirichlet Problem \dotfill 100
6.4. \ Interior and Boundary Regularity \dotfill 109
6.5. \ An Alternative Approach \dotfill 112
6.6. \ Non-Uniformly Elliptic Equations \dotfill 116
6.7. \ Other Boundary Conditions; the Oblique Derivative Problem \dotfill 120
6.8. \ Appendix 1: Interpolation Inequalities \dotfill 130
6.9. \ Appendix 2: Extension Lemmas \dotfill 136
Notes \dotfill 138
Problems \dotfill 141

Chapter 7. \ Sobolev Spaces \dotfill 144

7.1. \ $L^p$ Spaces \dotfill 145
7.2. \ Regularization and Approximation by Smooth Functions \dotfill 147
7.3. \ Weak Derivatives \dotfill 149
7.4. \ The Chain Rule \dotfill 151
7.5. \ The $W^{k,p}$ Spaces \dotfill 153
7.6. \ Density Theorems \dotfill 154
7.7. \ Imbedding Theorems \dotfill 155
7.8. \ Potential Estimates and Imbedding Theorems \dotfill 159
7.9. \ The Morrey and John-Nirenberg Estimates \dotfill 164
7.10. \ Compactness Results \dotfill 167
\end{quote}

% PDF page 11
\source{gilbarg-trudinger:page-0011}
\begin{quote}
Table of Contents

7.11. \ Difference Quotients \dotfill 168
7.12. \ Extension and Interpolation \dotfill 169
Notes \dotfill 173
Problems \dotfill 173

Chapter 8. \ Generalized Solutions and Regularity \dotfill 177

8.1. \ The Weak Maximum Principle \dotfill 179
8.2. \ Solvability of the Dirichlet Problem \dotfill 181
8.3. \ Differentiability of Weak Solutions \dotfill 183
8.4. \ Global Regularity \dotfill 186
8.5. \ Global Boundedness of Weak Solutions \dotfill 188
8.6. \ Local Properties of Weak Solutions \dotfill 194
8.7. \ The Strong Maximum Principle \dotfill 198
8.8. \ The Harnack Inequality \dotfill 199
8.9. \ Hölder Continuity \dotfill 200
8.10. \ Local Estimates at the Boundary \dotfill 202
8.11. \ Hölder Estimates for the First Derivatives \dotfill 209
8.12. \ The Eigenvalue Problem \dotfill 212
Notes \dotfill 214
Problems \dotfill 216

Chapter 9. \ Strong Solutions \dotfill 219

9.1. \ Maximum Principles for Strong Solutions \dotfill 220
9.2. \ $L^p$ Estimates: Preliminary Analysis \dotfill 225
9.3. \ The Marcinkiewicz Interpolation Theorem \dotfill 227
9.4. \ The Calderon-Zygmund Inequality \dotfill 230
9.5. \ $L^p$ Estimates \dotfill 235
9.6. \ The Dirichlet Problem \dotfill 241
9.7. \ A Local Maximum Principle \dotfill 244
9.8. \ Hölder and Harnack Estimates \dotfill 246
9.9. \ Local Estimates at the Boundary \dotfill 250
Notes \dotfill 254
Problems \dotfill 255

Part II. \ Quasilinear Equations

Chapter 10. \ Maximum and Comparison Principles \dotfill 259

10.1. \ The Comparison Principle \dotfill 263
10.2. \ Maximum Principles \dotfill 264
10.3. \ A Counterexample \dotfill 267
10.4. \ Comparison Principles for Divergence Form Operators \dotfill 268
10.5. \ Maximum Principles for Divergence Form Operators \dotfill 271
Notes \dotfill 277
Problems \dotfill 277
\end{quote}

% PDF page 12
\source{gilbarg-trudinger:page-0012}

\begin{quote}
\hfill \textit{Table of Contents}

\vspace{1em}

\noindent \textbf{Chapter 11.}  Topological Fixed Point Theorems and Their Application \hfill 279

\noindent \hspace*{1.5em} 11.1. \ \ The Schauder Fixed Point Theorem \hfill 279

\noindent \hspace*{1.5em} 11.2. \ \ The Leray-Schauder Theorem: a Special Case \hfill 280

\noindent \hspace*{1.5em} 11.3. \ \ An Application \hfill 282

\noindent \hspace*{1.5em} 11.4. \ \ The Leray-Schauder Fixed Point Theorem \hfill 286

\noindent \hspace*{1.5em} 11.5. \ \ Variational Problems \hfill 288

\noindent Notes \hfill 293

\vspace{1em}

\noindent \textbf{Chapter 12.}  Equations in Two Variables \hfill 294

\noindent \hspace*{1.5em} 12.1. \ \ Quasiconformal Mappings \hfill 294

\noindent \hspace*{1.5em} 12.2. \ \ Hölder Gradient Estimates for Linear Equations \hfill 300

\noindent \hspace*{1.5em} 12.3. \ \ The Dirichlet Problem for Uniformly Elliptic Equations \hfill 304

\noindent \hspace*{1.5em} 12.4. \ \ Non-Uniformly Elliptic Equations \hfill 309

\noindent Notes \hfill 315

\noindent Problems \hfill 317

\vspace{1em}

\noindent \textbf{Chapter 13.}  Hölder Estimates for the Gradient \hfill 319

\noindent \hspace*{1.5em} 13.1. \ \ Equations of Divergence Form \hfill 319

\noindent \hspace*{1.5em} 13.2. \ \ Equations in Two Variables \hfill 323

\noindent \hspace*{1.5em} 13.3. \ \ Equations of General Form; the Interior Estimate \hfill 324

\noindent \hspace*{1.5em} 13.4. \ \ Equations of General Form; the Boundary Estimate \hfill 328

\noindent \hspace*{1.5em} 13.5. \ \ Application to the Dirichlet Problem \hfill 331

\noindent Notes \hfill 332

\vspace{1em}

\noindent \textbf{Chapter 14.}  Boundary Gradient Estimates \hfill 333

\noindent \hspace*{1.5em} 14.1. \ \ General Domains \hfill 335

\noindent \hspace*{1.5em} 14.2. \ \ Convex Domains \hfill 337

\noindent \hspace*{1.5em} 14.3. \ \ Boundary Curvature Conditions \hfill 341

\noindent \hspace*{1.5em} 14.4. \ \ Non-Existence Results \hfill 347

\noindent \hspace*{1.5em} 14.5. \ \ Continuity Estimates \hfill 353

\noindent \hspace*{1.5em} 14.6. \ \ Appendix: Boundary Curvatures and the Distance Function \hfill 354

\noindent Notes \hfill 357

\noindent Problems \hfill 358

\vspace{1em}

\noindent \textbf{Chapter 15.}  Global and Interior Gradient Bounds \hfill 359

\noindent \hspace*{1.5em} 15.1. \ \ A Maximum Principle for the Gradient \hfill 359

\noindent \hspace*{1.5em} 15.2. \ \ The General Case \hfill 362

\noindent \hspace*{1.5em} 15.3. \ \ Interior Gradient Bounds \hfill 369

\noindent \hspace*{1.5em} 15.4. \ \ Equations in Divergence Form \hfill 373

\noindent \hspace*{1.5em} 15.5. \ \ Selected Existence Theorems \hfill 380

\noindent \hspace*{1.5em} 15.6. \ \ Existence Theorems for Continuous Boundary Values \hfill 384

\noindent Notes \hfill 385

\noindent Problems \hfill 386
\end{quote}

% PDF page 13
\source{gilbarg-trudinger:page-0013}
\nobreak
\noindent Table of Contents \hfill xiii

\medskip

\noindent Chapter 16.  Equations of Mean Curvature Type \dotfill 388

\noindent\hspace*{1.5em}16.1.  Hypersurfaces in $\mathbb R^{n+1}$ \dotfill 388

\noindent\hspace*{1.5em}16.2.  Interior Gradient Bounds \dotfill 401

\noindent\hspace*{1.5em}16.3.  Application to the Dirichlet Problem \dotfill 407

\noindent\hspace*{1.5em}16.4.  Equations in Two Independent Variables \dotfill 410

\noindent\hspace*{1.5em}16.5.  Quasiconformal Mappings \dotfill 413

\noindent\hspace*{1.5em}16.6.  Graphs with Quasiconformal Gauss Map \dotfill 423

\noindent\hspace*{1.5em}16.7.  Applications to Equations of Mean Curvature Type \dotfill 429

\noindent\hspace*{1.5em}16.8.  Appendix: Elliptic Parametric Functionals \dotfill 434

\noindent Notes \dotfill 437

\noindent Problems \dotfill 438

\bigskip

\noindent Chapter 17.  Fully Nonlinear Equations \dotfill 441

\noindent\hspace*{1.5em}17.1.  Maximum and Comparison Principles \dotfill 443

\noindent\hspace*{1.5em}17.2.  The Method of Continuity \dotfill 446

\noindent\hspace*{1.5em}17.3.  Equations in Two Variables \dotfill 450

\noindent\hspace*{1.5em}17.4.  Hölder Estimates for Second Derivatives \dotfill 453

\noindent\hspace*{1.5em}17.5.  Dirichlet Problem for Uniformly Elliptic Equations \dotfill 463

\noindent\hspace*{1.5em}17.6.  Second Derivative Estimates for Equations of
\par\noindent\hspace*{3.0em}Monge-Ampère Type \dotfill 467

\noindent\hspace*{1.5em}17.7.  Dirichlet Problem for Equations of Monge-Ampère Type \dotfill 471

\noindent\hspace*{1.5em}17.8.  Global Second Derivative Hölder Estimates \dotfill 476

\noindent\hspace*{1.5em}17.9.  Nonlinear Boundary Value Problems \dotfill 481

\noindent Notes \dotfill 486

\noindent Problems \dotfill 488

\bigskip

\noindent Bibliography \dotfill 491

\bigskip

\noindent Epilogue \dotfill 507

\bigskip

\noindent Subject Index \dotfill 511

\bigskip

\noindent Notation Index \dotfill 516
